\documentclass[11pt,a4paper]{article}
\usepackage[utf8]{inputenc}
\usepackage[T1]{fontenc}
\usepackage[margin=1in]{geometry}
\usepackage{booktabs}
\usepackage{tabularx}
\usepackage{array}

\newcolumntype{L}[1]{>{\raggedright\arraybackslash}p{#1}}
\newcolumntype{Y}{>{\raggedright\arraybackslash}X}

\title{Common NDVI Bins and Their Interpretation}
\date{}

\begin{document}

\maketitle

\section*{Standard Binned NDVI Classification}

While exact thresholds vary slightly depending on the satellite sensor, atmospheric correction, and local biome, remote-sensing classifications generally interpret Normalized Difference Vegetation Index (NDVI) ranges as follows:

\begin{table}[h!]
\centering
\renewcommand{\arraystretch}{1.35}
\begin{tabularx}{\textwidth}{L{2.4cm} L{4.2cm} Y}
\toprule
\textbf{NDVI Range (Bin)} & \textbf{Surface / Vegetation Class} & \textbf{Interpretation} \\
\midrule
$-1.0$ to $0.0$ & Water, Snow, Ice, or Clouds & Deep water typically falls near $-1.0$, while shallow or turbid water, snow, and thick clouds range between $-0.2$ and $0.0$. \\
$0.0$ to $0.1$ & Barren Rock, Sand, or Urban Surfaces & Impervious surfaces (concrete, asphalt), exposed bedrock, dry sand, or completely dead/dormant ground cover. \\
$0.1$ to $0.2$ & Bare Soil or Very Sparse Cover & Open soil with minimal emerging seedlings, dry desert scrub, or harvested agricultural fields. \\
$0.2$ to $0.4$ & Sparse to Moderate Vegetation & Grasslands, shrublands, senescing (aging) crops, or vegetation experiencing moderate drought or nutrient stress. \\
$0.4$ to $0.6$ & Moderate to Healthy Vegetation & Active agricultural crops, open woodlands, healthy pastures, and temperate mixed cover. \\
$0.6$ to $1.0$ & Dense, Vigorous Vegetation & Closed-canopy forests (e.g., tropical rainforests, peak-season deciduous forests) and peak-growth irrigated crops. \\
\bottomrule
\end{tabularx}
\caption{Standard interpretation of binned NDVI values across land-cover types.}
\end{table}

\end{document}
